\documentclass[10pt,a4paper]{article}
\usepackage[a4paper,margin=14mm,top=12mm,bottom=14mm]{geometry}
\usepackage{fontspec}
\setmainfont{Liberation Sans}
\usepackage{polyglossia}
\setdefaultlanguage{french}
\usepackage{amsmath}
\usepackage{xcolor}
\usepackage{graphicx}
\usepackage[most]{tcolorbox}
\usepackage{tabularx,booktabs,array,multirow}
\usepackage{enumitem}
\usepackage{fancyhdr}
\usepackage{lastpage}
\usepackage[hidelinks]{hyperref}
\definecolor{uimm}{HTML}{005677}
\definecolor{cyan}{HTML}{13A2DD}
\definecolor{fond}{HTML}{E6F2FA}
\definecolor{ok}{HTML}{2E7D32}
\definecolor{ko}{HTML}{B3261E}
\definecolor{violet}{HTML}{7B1FA2}
\setlength{\parindent}{0pt}
\setlength{\parskip}{3pt}
\setlist{nosep}
\setlist[itemize]{leftmargin=13pt,label=\textcolor{cyan}{\textbullet}}
\newcommand{\titre}[1]{\vspace{5pt}{\color{uimm}\large\bfseries #1}\par\vspace{-3pt}{\color{cyan}\rule{\linewidth}{1.2pt}}\par\vspace{0pt}}
\tcbset{boite/.style={colback=fond,colframe=uimm,boxrule=0.6pt,arc=2pt,left=6pt,right=6pt,top=4pt,bottom=4pt},
        defi/.style={colback=violet!6,colframe=violet,boxrule=0.8pt,arc=2pt,left=6pt,right=6pt,top=4pt,bottom=4pt,
                     title={\bfseries Défi},coltitle=white,colbacktitle=violet,fonttitle=\small}}
\newcommand{\bandeau}[3]{%
\begin{tcolorbox}[colback=uimm,colframe=uimm,arc=2pt,left=4pt,right=8pt,top=4pt,bottom=4pt]
\begin{minipage}[c]{18mm}\setlength{\fboxsep}{1.2pt}\fcolorbox{white}{white}{\includegraphics[height=16.5mm]{logo_uimm.jpg}}\end{minipage}\hspace{4mm}%
\begin{minipage}[c]{\dimexpr\linewidth-22mm\relax}\color{white}
{\small #1}\\[1pt]{\LARGE\bfseries #2}\\[2pt]{\large #3}\hfill{\small CQPM Technicien qualité}
\end{minipage}
\end{tcolorbox}}
\newcommand{\reponse}[1][1]{\par\vspace{2pt}{\color{gray!60}\foreach \i in {1,...,#1}{\rule{\linewidth}{0.3pt}\par\vspace{9pt}}}}
\usepackage{pgffor}
\pagestyle{fancy}\fancyhf{}
\renewcommand{\headrulewidth}{0pt}
\newcommand{\pied}[1]{\fancyfoot[L]{\footnotesize\color{uimm}#1}\fancyfoot[R]{\footnotesize\color{uimm}Page \thepage{} sur \pageref{LastPage}}}

\usepackage{tikz}
\usepackage{longtable}
\renewcommand{\bandeau}[3]{%
\begin{tcolorbox}[colback=uimm,colframe=uimm,arc=2pt,left=4pt,right=8pt,top=4pt,bottom=4pt]
\begin{minipage}[c]{18mm}\setlength{\fboxsep}{1.2pt}\fcolorbox{white}{white}{\includegraphics[height=16.5mm]{logo_uimm.jpg}}\end{minipage}\hspace{4mm}%
\begin{minipage}[c]{\dimexpr\linewidth-22mm\relax}\color{white}
{\small #1}\\[1pt]{\LARGE\bfseries #2}\\[2pt]{\large #3}
\end{minipage}
\end{tcolorbox}}
% en-tête d'une fiche : code, titre, question, bloc, durée, organisation, niveau
\newcommand{\entete}[7]{%
\begin{tcolorbox}[colback=uimm,colframe=uimm,arc=2pt,left=4pt,right=6pt,top=4pt,bottom=4pt]
\begin{minipage}[c]{18mm}\setlength{\fboxsep}{1.2pt}\fcolorbox{white}{white}{\includegraphics[height=16.5mm]{logo_uimm.jpg}}\end{minipage}\hspace{3mm}%
\begin{minipage}[c]{\dimexpr\linewidth-58mm\relax}\color{white}\raggedright
{\small Base SPC · #4}\\[1pt]{\Large\bfseries #1 · #2}\\[2pt]{\itshape #3}
\end{minipage}\hfill
\begin{minipage}[c]{33mm}\color{white}\small\raggedleft
Durée : \textbf{#5}\\ \textbf{#6}\\ Niveau #7 sur 3
\end{minipage}
\end{tcolorbox}}
\newcommand{\tagE}{{\setlength{\fboxsep}{1.5pt}\colorbox{cyan}{\color{white}\scriptsize\bfseries Ellistat}}}
\newcommand{\tagW}{{\setlength{\fboxsep}{1.5pt}\colorbox{uimm}{\color{white}\scriptsize\bfseries Page web}}}
\tcbset{pred/.style={colback=orange!6,colframe=orange!80!black,boxrule=0.6pt,arc=2pt,left=6pt,right=6pt,top=3pt,bottom=1pt,
                    title={\bfseries Je prédis, avant de commencer},coltitle=white,colbacktitle=orange!80!black,fonttitle=\small},
        retenir/.style={colback=fond,colframe=uimm,boxrule=0.6pt,arc=2pt,left=6pt,right=6pt,top=3pt,bottom=3pt,
                    title={\bfseries À retenir},coltitle=white,colbacktitle=uimm,fonttitle=\small},
        plus/.style={colback=violet!6,colframe=violet,boxrule=0.6pt,arc=2pt,left=6pt,right=6pt,top=3pt,bottom=3pt,
                    title={\bfseries Pour aller plus loin},coltitle=white,colbacktitle=violet,fonttitle=\small},
        act/.style={colback=uimm!10,colframe=uimm,boxrule=0.5pt,arc=2pt,left=5pt,right=5pt,top=2pt,bottom=2pt,before skip=7pt,after skip=3pt},
        verif/.style={colback=orange!6,colframe=orange!80!black,boxrule=0.6pt,arc=2pt,left=6pt,right=6pt,top=3pt,bottom=3pt,
                    title={\bfseries À vérifier sur poste avant la première séance},coltitle=white,colbacktitle=orange!80!black,fonttitle=\small}}
\newcommand{\fichier}[1]{\texttt{\small #1}}
\newcommand{\menu}[1]{\textsf{\textbf{#1}}}
\newcommand{\champ}[1]{\rule{#1}{0.3pt}}

\usepackage{amssymb}
\newcommand{\cb}{$\square$\,}
\newcommand{\etape}[2]{\par\vspace{9pt}{\setlength{\fboxsep}{3pt}\colorbox{uimm}{\color{white}\bfseries\large\ #1\ }}\hspace{4pt}{\color{uimm}\large\bfseries #2}\par\vspace{2pt}}
\newcommand{\motnouveau}[1]{\begin{tcolorbox}[colback=violet!6,colframe=violet,boxrule=0.6pt,arc=2pt,left=6pt,right=6pt,top=3pt,bottom=3pt,title={\bfseries Mot nouveau},coltitle=white,colbacktitle=violet,fonttitle=\small]#1\end{tcolorbox}}
\newcommand{\chiffre}[1]{\begin{tcolorbox}[colback=orange!8,colframe=orange!80!black,boxrule=1pt,arc=3pt,left=8pt,right=8pt,top=5pt,bottom=5pt]#1\end{tcolorbox}}

\pied{Journée du 25/11/2026 · Trames terrain · repères du formateur}
\begin{document}
\thispagestyle{fancy}
\bandeau{Journée du 25/11/2026 · 4 stagiaires}{Trames terrain : repères du formateur}{Cinq TP d'une page, à imprimer en 4 exemplaires}
\titre{Ordre conseillé et durées}
\renewcommand{\arraystretch}{1.35}
\begin{tabularx}{\linewidth}{|c|>{\raggedright\arraybackslash}p{40mm}|c|>{\raggedright\arraybackslash}p{37mm}|>{\raggedright\arraybackslash}X|}\hline
\textbf{TP} & \textbf{Titre (fiche de la base)} & \textbf{Durée} & \textbf{Matériel} & \textbf{Chiffre lu, valeurs du jeu de secours}\\\hline
1 & Mon micromètre est-il juste ? (M3) & 45 min & micromètre 0-25 et cale 10 mm par binôme & Cgk. Secours : 6,02 pour $T = 0{,}15$ (formules du livret Bosch).\\\hline
2 & Jusqu'où peut-il aller ? (M7) & 15 min & aucun & Tolerance min. Secours : 0,036 mm, donc capable à 0,05 et pas à 0,02.\\\hline
3 & Pied à coulisse contre micromètre (M4) & 30 min & pied à coulisse au 1/50 & Cgk du pied à coulisse. Secours : 0,45 : pas capable à 0,15 ; résolution 0,02 permet $T \geq 0{,}4$ mm.\\\hline
4 & Avec ou sans cliquet (M5) & 30 min & mêmes micromètre et cale & Écart type. Secours : 0,0005 avec cliquet contre 0,0023 sans.\\\hline
5 & La cale dans la main (M6) & 25 min & cale longue, chronomètre & Allongement. Secours : environ +3 µm, retour en 3 minutes.\\\hline
\end{tabularx}
\par\smallskip
Le TP 2 vient juste après le TP 1 : il réutilise l'étude sans nouvelle mesure. L'ordre de la journée, les durées et les réserves sont dans la trame de 8 h 20.

\titre{Ce qui reste à vérifier sur ton compte avant la séance}
\begin{itemize}
\item \textbf{Saisie d'une justesse sur une cale de 10 mm} : Ellistat attend une valeur vraie, LSI et LSS. Nous centrons LSI et LSS sur la cale (largeur 0,15). Seule la largeur est confirmée par les formules ; l'effet de la position de LSI et LSS n'a pas été testé.
\item \textbf{Pied à coulisse} : vérifier que les valeurs 10,00 et 10,02 sont acceptées et ce que l'écran affiche. Non vu.
\item \textbf{Onglet Détail avec plusieurs colonnes} (TP 4 et 5) : vu sur une seule colonne. Le clic sur le nom d'une colonne dans la liste « Échantillons » est supposé.
\item \textbf{Seuil 1,33} : c'est celui du livret Bosch. Les seuils de couleur d'Ellistat n'ont pas été lus.
\item \textbf{Valeur des cales} : écart de la valeur vraie à lire sur le certificat des cales ; à défaut, la valeur nominale est utilisée et le dire.
\item \textbf{TP 5} : la dilatation de 11,5 µm par mètre et par degré est une valeur usuelle de l'acier, à confirmer sur la documentation de tes cales. Avec une cale de 20 mm, l'effet n'est que d'environ 0,23 µm par degré : la cale de 30 mm (micromètre 25-50, un seul) montre mieux, avec une rotation.
\item \textbf{Frappe de 25 valeurs} : prévoir un secrétaire par binôme. Le collage d'une liste dans la fenêtre « Justesse » n'a pas été testé.
\end{itemize}
\titre{Comment s'en servir}
Les trames sont des versions courtes de M3, M7, M4, M5 et M6. Les fiches complètes (questions, réponses, « pour aller plus loin ») restent dans le livret et le guide formateur. Chaque trame a une seule lecture dans Ellistat et au plus un mot nouveau. Faire annoncer la prédiction avant de mesurer, et faire dire la décision à voix haute en fin de TP.
\end{document}
